% =========================================================
% DOCUMENTCLASS
% =========================================================
\documentclass[12pt,oneside]{book}

% =========================================================
% METADATOS
% =========================================================
\newcommand{\universidad}{Universidad de Buenos Aires (UBA)}
\newcommand{\materia}{Derecho Procesal Civil y Comercial}
\newcommand{\titulo}{Diligencias Preliminares y la Demanda}
\newcommand{\autor}{Isaac Edgar Camacho Ocampo}
\newcommand{\anio}{2026}

% =========================================================
% PAQUETES
% =========================================================
\usepackage[utf8]{inputenc}
\usepackage[T1]{fontenc}
\usepackage[spanish,es-nodecimaldot]{babel}

\usepackage[
    top=2.5cm,
    bottom=2.5cm,
    left=3cm,
    right=2.5cm,
    headheight=15pt
]{geometry}

\usepackage{amsmath}
\usepackage{amssymb}
\usepackage{mathtools}

\usepackage{graphicx}
\usepackage{float}
\usepackage{caption}
\usepackage{subcaption}

\usepackage{booktabs}
\usepackage{array}
\usepackage{longtable}
\usepackage{tabularx}

\usepackage{enumitem}

\usepackage{xcolor}
\usepackage[most]{tcolorbox}

\usepackage{fancyhdr}
\usepackage{ragged2e}

\usepackage[
    colorlinks=true,
    linkcolor=blue!50!black,
    citecolor=green!40!black,
    urlcolor=blue!60!black,
    pdftitle={\titulo},
    pdfauthor={\autor},
    pdfsubject={Apuntes universitarios},
    bookmarks=true
]{hyperref}

% =========================================================
% CONFIGURACIÓN
% =========================================================
\graphicspath{
    {./img/}
    {./graficos/}
    {./figuras/}
}

\setcounter{tocdepth}{2}

\setlength{\parindent}{0pt}
\setlength{\parskip}{0.5em}

\setlist[itemize]{
    topsep=0.3em,
    itemsep=0.2em
}

\setlist[enumerate]{
    topsep=0.3em,
    itemsep=0.2em
}

\clubpenalty=10000
\widowpenalty=10000

% =========================================================
% COMANDOS
% =========================================================
\newcommand{\abs}[1]{\left|#1\right|}
\newcommand{\norm}[1]{\left\lVert#1\right\rVert}
\newcommand{\vect}[1]{\mathbf{#1}}
\newcommand{\deriv}[2]{\frac{d#1}{d#2}}
\newcommand{\pderiv}[2]{\frac{\partial#1}{\partial#2}}

% =========================================================
% ENTORNOS / CAJAS
% =========================================================
\newtcolorbox{definition}[1]{
    enhanced,
    breakable,
    title={Definición: #1},
    fonttitle=\bfseries,
    colback=blue!3,
    colframe=blue!50!black,
    boxrule=0.8pt,
    arc=2mm
}

\newtcolorbox{important}{
    enhanced,
    breakable,
    title={Importante},
    fonttitle=\bfseries,
    colback=yellow!8,
    colframe=orange!70!black,
    boxrule=0.8pt,
    arc=2mm
}

\newtcolorbox{example}[1]{
    enhanced,
    breakable,
    title={Ejemplo: #1},
    fonttitle=\bfseries,
    colback=green!4,
    colframe=green!40!black,
    boxrule=0.8pt,
    arc=2mm
}

\newtcolorbox{formula}[1]{
    enhanced,
    breakable,
    title={Fórmula / Regla: #1},
    fonttitle=\bfseries,
    colback=purple!4,
    colframe=purple!50!black,
    boxrule=0.8pt,
    arc=2mm
}

\newtcolorbox{exercise}[1]{
    enhanced,
    breakable,
    title={Ejercicio: #1},
    fonttitle=\bfseries,
    colback=gray!5,
    colframe=gray!60!black,
    boxrule=0.8pt,
    arc=2mm
}

\newtcolorbox{note}{
    enhanced,
    breakable,
    title={Nota},
    fonttitle=\bfseries,
    colback=white,
    colframe=gray!50,
    boxrule=0.6pt,
    arc=2mm
}

% =========================================================
% CONFIGURACIÓN FINAL (encabezados y pies)
% =========================================================
\pagestyle{fancy}
\fancyhf{}

\fancyhead[L]{\nouppercase{\leftmark}}
\fancyhead[R]{\materia}

\fancyfoot[C]{\thepage}

\renewcommand{\headrulewidth}{0.4pt}
\renewcommand{\footrulewidth}{0pt}

\fancypagestyle{plain}{
    \fancyhf{}
    \fancyfoot[C]{\thepage}
    \renewcommand{\headrulewidth}{0pt}
}

% =========================================================
% DOCUMENT
% =========================================================
\begin{document}

% ---------------------------------------------------------
% PORTADA
% ---------------------------------------------------------
\begin{titlepage}

\thispagestyle{empty}

\begin{center}

\vspace*{1cm}

{\large \textsc{\universidad}}

\vspace{3cm}

{\Huge \textbf{\materia}}

\vspace{1cm}

{\LARGE \titulo}

\vspace{0.8cm}

\textit{Comprender los conceptos es el primer paso para convertir el estudio en conocimiento.}

\vspace{1.2cm}

\rule{0.70\textwidth}{0.5pt}

\vspace{0.5cm}

{\large Apuntes de estudio}

\vspace{0.5cm}

\rule{0.70\textwidth}{0.5pt}

\vfill

{\large \autor}

\vspace{0.3cm}

{\large \anio}

\end{center}

\vfill

\begin{flushleft}
\textit{Este es un modesto aporte a los estudiantes de ``Derecho Procesal Civil y Comercial'' no pretende sustituir el material oficial de la materia.}
\end{flushleft}

\end{titlepage}

% ---------------------------------------------------------
% ÍNDICE
% ---------------------------------------------------------
\tableofcontents

% =========================================================
% CONTENIDO
% =========================================================

\chapter{Introducción: la demanda y las diligencias preliminares}

\section{La demanda como acto procesal. Concepto inicial}

La demanda debe pensarse como un acto procesal mediante el cual se formula una petición dotada de determinadas formalidades, cuya finalidad es dar inicio al proceso judicial en el que habrá de ventilarse un conflicto suscitado entre dos o más personas. Para resolver dicho conflicto es necesario recurrir al derecho de acción, entendido como una potencialidad que corresponde a toda persona por el solo hecho de serlo. A través de la concreción de ese derecho de acción surge una pretensión, que es concretamente aquello que se le pide al juez que otorgue mediante el dictado de la sentencia. Esa pretensión debe incluirse dentro de una demanda, que constituye su continente, y a través de ella se efectúa la petición relativa a que se conceda dicha pretensión.

\begin{definition}{Demanda}
La demanda es un acto de petición que contiene una pretensión, integrada por los elementos ya estudiados con anterioridad.
\end{definition}

Los elementos de la pretensión son los siguientes:

\begin{itemize}
    \item \textbf{Elemento objetivo}: comprende el objeto inmediato y el objeto mediato, y la causa, que se vincula con los hechos en los que se fundamenta la pretensión.
    \item \textbf{Elementos subjetivos}: comprenden la parte que reclama la satisfacción de la pretensión (normalmente, la parte actora) y la parte frente a quien se reclama dicha satisfacción (el demandado). En un plano superior, ordenado respecto de las partes, se ubica el juez como director del proceso.
\end{itemize}

\section{Concepto y ubicación de las diligencias preliminares en el Código}

Las medidas preparatorias sirven para que quien vaya a ser parte en un proceso —como actor o como demandado— requiera, cuando no puede hacerlo extrajudicialmente, que se lleven a cabo ciertas diligencias destinadas a preparar adecuadamente la demanda o a dotar de mayor eficacia al proceso judicial.

En cuanto a su ubicación, el Código Procesal aborda el tema en la parte especial, dentro del Libro Segundo —que trata los procesos de conocimiento—, en el Título Primero de disposiciones generales. Luego de que el Capítulo Primero se ocupa de las clases de proceso, el Capítulo Segundo regula las diligencias preliminares, en los artículos 323 a 329.

\begin{important}
El Código trata, dentro de las diligencias preliminares, dos institutos que se diferencian perfectamente entre sí: las \textbf{diligencias preparatorias} y la \textbf{prueba anticipada}. Tienen distinta naturaleza jurídica y apuntan a distintas cuestiones.
\end{important}

\chapter{Medidas preparatorias}

\section{Concepto de las medidas preparatorias}

Las medidas preparatorias son una serie de posibilidades que le asisten a quien vaya a demandar —es decir, al actor que piensa iniciar una demanda— o a quien, con fundamento, según indica el Código, prevé que será demandado, con el fin de obtener una serie de datos o elementos sin los cuales no podría iniciarse válidamente la demanda.

Cuando la parte no puede obtener extrajudicialmente esos datos o elementos, el Código permite, a través de las diligencias preparatorias, pedirle al juez que tome intervención para ordenarlas y hacerlas efectivas, con distintas consecuencias para el caso de que quien vaya a ser demandado, por ejemplo, no cumpla con la manda judicial.

\section{Procesos en los que se pueden solicitar}

Las medidas preparatorias están previstas dentro de los artículos correspondientes a los procesos de conocimiento, es decir, el proceso ordinario y el sumarísimo —dado que el proceso sumario fue derogado—. La doctrina entiende, además, que resultan aplicables a los procesos especiales que tramitan bajo las formas del proceso ordinario o sumarísimo.

También corresponde una medida en el proceso sucesorio, contemplada en uno de los once incisos del artículo 323, que le permite a quien se crea heredero o legatario solicitar que se exhiba un testamento a los efectos de ejercer sus derechos.

\section{Caracteres de las medidas preparatorias}

\begin{enumerate}[label=\alph*)]
    \item \textbf{Son excepcionales o taxativas}: solo pueden solicitarse las medidas contempladas en los once incisos del artículo 323. Sin perjuicio de ello, cabe señalar que la jurisprudencia ha ido morigerando esta taxatividad, admitiendo medidas preparatorias no contempladas específicamente cuando su falta de ordenamiento, a pedido de parte, terminaría perjudicando la eficacia de una demanda o imposibilitando su interposición, conculcando el derecho que se perseguiría hacer efectivo en ese proceso.
    \item \textbf{No provocan la apertura de la instancia}: el acto procesal típico para abrir la instancia judicial es la demanda. La instancia —una de las etapas visibles del proceso— se abre con la presentación de la demanda, y no con el pedido de una medida preparatoria del artículo 323.
    \item \textbf{Interrumpen el curso de la prescripción}: la interposición de una diligencia o medida preparatoria interrumpe el curso de la prescripción.
    \item \textbf{Determinación del juez competente}: será competente el juez que corresponda intervenir en el proceso que habrá de promoverse, dado que estas medidas tienen por finalidad la preparación precisa y eficaz de una demanda que provocará la apertura de un proceso.
    \item \textbf{Oportunidad}: pueden pedirse antes de interpuesta la demanda —lo habitual— o bien luego de interpuesta, pero siempre antes de la notificación del traslado al demandado, momento en el cual queda trabada la litis.
\end{enumerate}

\section{Contenido de las medidas preparatorias: incisos del artículo 323}

Para determinar en qué consisten las medidas preparatorias que pueden solicitarse, corresponde remitirse a los incisos del artículo 323.

\begin{example}{Inciso primero del artículo 323}
El inciso primero permite pedir la declaración jurada respecto de algún dato que hace a la personalidad de quien se pretende demandar, y sin el cual no podría entrarse en juicio. Estos datos o hechos relativos a la personalidad del futuro demandado no deben apuntar al fundamento de la pretensión: no es este el momento para acreditar, a través de las medidas preparatorias, hechos que deberán discutirse una vez instalado el proceso judicial. Pueden recabarse, así, datos relativos a la capacidad de una persona, a su domicilio o a su legitimación procesal.
\end{example}

% TODO: contenido incompleto o ambiguo en las notas originales — el apunte solo desarrolla en detalle el inciso primero del artículo 323; los restantes incisos no fueron desarrollados en la clase.

\section{Trámite de las medidas preparatorias}

La parte que requiere una medida preparatoria se presenta mediante un escrito, que debe reunir todas las formalidades correspondientes a las postulaciones de las partes en el proceso. Dichos escritos deben confeccionarse conforme a lo establecido por el artículo 115 y siguientes del Código Procesal y, en su caso, por el Reglamento de la Justicia Nacional, artículos 45, 46 y 47.

\begin{note}
Las medidas preparatorias se dictan \textit{sin sustanciación}. Esta expresión no significa que carezcan de trámite —pues todo trámite lo tiene—, sino que se resuelven sin controvertir el trámite con la otra parte. Sustanciar significa transitar una cuestión por los carriles adecuados para que quede en estado de ser resuelta, lo cual no siempre requiere darle traslado a la contraria para que tenga la posibilidad de controvertir o no la posición contraria.
\end{note}

El juez hace lugar o no a la medida preparatoria en función de los fundamentos que da el actor —si es quien la solicita— o, en su caso, el demandado, si con fundamento prevé que será demandado.

\section{Incumplimiento y sanciones (artículo 329)}

El artículo 329 indica distintas consecuencias ante el eventual incumplimiento por parte del requerido. Existen sanciones procesales que van desde la multa hasta la aplicación de astreintes, e incluso la presunción en contra de quien adopta una actitud omisiva respecto del requerimiento efectuado por el juez.

Según cuál sea la medida de que se trate —por ejemplo, si se pide como medida preparatoria que se exhiban cosas o documentos y no se exhiben—, puede solicitarse incluso el secuestro, con la eventual orden de allanamiento para hacerlo efectivo, o la aplicación de astreintes. Existe, en definitiva, un marco de herramientas de las que puede echarse mano.

\section{Caducidad de las medidas preparatorias}

\begin{important}
Una vez hechas efectivas, las medidas preparatorias caducan a los treinta días de que se tuvo por cumplida la resolución que las ordena, si no se promueve la demanda. Se exceptúan las contempladas en los incisos 9, 10 y 11 del artículo 323 —referidas al reconocimiento de la obligación de rendir cuentas, al reconocimiento de mercaderías y a la mensura—, que no caducan.
\end{important}

\chapter{Prueba anticipada}

\section{Concepto y naturaleza jurídica}

El Código Procesal contempla la prueba anticipada en el artículo 326, en cuatro incisos —el último de ellos agregado por la reforma de la ley 25.488— y con una última oración que alude a la prueba confesional como prueba anticipada.

\begin{definition}{Prueba}
La prueba consiste en la demostración, dentro del proceso, de la ocurrencia de los hechos que deben acreditarse en el proceso judicial.
\end{definition}

La anticipación a la que alude la prueba anticipada se refiere al momento procesal establecido por el Código para que la prueba se produzca, que es la etapa probatoria. En los escritos constitutivos del proceso —demanda, contestación de demanda, reconvención, contestación de reconvención, excepciones previas y contestación de excepciones previas— las partes ofrecen la prueba, es decir, indican de qué medios se quieren valer; pero su producción debe tener lugar en la etapa probatoria, una vez abierta la causa a prueba, lo cual se decide en la audiencia preliminar cuando existan hechos respecto de los cuales no haya conformidad.

\section{Naturaleza cautelar de la prueba anticipada}

La necesidad de evitar los perjuicios que pueden ocasionar las distintas instancias que requiere el proceso, con los tiempos consiguientes, exige alguna actuación del juez tendiente a conservar la situación probatoria y a evitar que el daño se produzca. En materia probatoria sucede lo mismo: si resulta necesario que declare un testigo de avanzada edad, próximo a ausentarse del país, puede actuarse de modo conservativo, permitiendo que ese testigo declare antes de la etapa probatoria.

\begin{note}
La palabra ``cautela'' proviene del latín y se vincula con ser cauto, obrar con cuidado y previsión, y ser precavido frente a una situación que puede ocasionar un perjuicio. La prueba anticipada participa de esta naturaleza cautelar porque el perjuicio que se busca evitar es la pérdida de la prueba por el paso del tiempo. Se la incluye dentro de los sistemas cautelares porque, si bien tiene connotaciones diferentes de otras medidas cautelares, comparte con ellas la finalidad de actuar ante el paso del tiempo y el peligro en la demora, para conservar una situación que, de no ser preservada, pondría a una de las partes en inferioridad de condiciones respecto de la otra.
\end{note}

\section{Ejemplos de prueba anticipada}

% TODO: contenido incompleto o ambiguo en las notas originales — el apunte anuncia este punto como apartado propio pero no desarrolla ejemplos concretos de prueba anticipada.

\section{Bilateralidad: la citación de la contraparte}

\begin{important}
A diferencia de lo que ocurre con las medidas preparatorias, en la prueba anticipada debe citarse a la contraparte para que controle la producción de la prueba. Toda prueba que se incorpore al proceso —por el medio de prueba anticipada que sea— debe permitir el control de la contraria; de lo contrario, esa prueba es nula. Este principio se vincula con la bilateralidad, la audiencia y el derecho de defensa.
\end{important}

Aun en los casos en que no sea conveniente citar a la contraparte —porque ello podría desnaturalizar la medida probatoria—, se requiere la citación del defensor oficial, quien debe intervenir y controlar la regularidad del acto cuando la medida probatoria no pueda llevarse adelante con notificación previa a la contraria sin que ello la perjudique.

\section{Apelabilidad de la denegatoria}

Denegada la medida de prueba anticipada, dicha denegatoria puede ser apelada. Se trata de una excepción al régimen general establecido por el artículo 379 del Código, según el cual toda medida vinculada con la denegación, sustanciación o producción de prueba es inapelable, sin perjuicio del replanteo que puede efectuarse en cámara cuando tramite el recurso de apelación.

% ---------------------------------------------------------
% CORNELL 1 — DILIGENCIAS PRELIMINARES
% ---------------------------------------------------------
\section*{Repaso: diligencias preliminares}
\addcontentsline{toc}{section}{Repaso: diligencias preliminares}

\begin{table}[H]
\centering
\small
\renewcommand{\arraystretch}{1.05}
\setlength{\parskip}{0.2em}
\begin{tabularx}{\textwidth}{
    >{\raggedright\arraybackslash}p{0.30\textwidth}
    >{\raggedright\arraybackslash}X
}
\toprule
\textbf{Preguntas / palabras clave} &
\textbf{Notas / respuestas} \\
\midrule

\textbf{¿Qué dos institutos comprenden las diligencias preliminares?} &
Las diligencias preparatorias y la prueba anticipada, que tienen distinta naturaleza jurídica y apuntan a distintas cuestiones. \\

\textbf{¿Dónde se regulan las diligencias preliminares en el Código?} &
En el Libro Segundo, Título Primero, Capítulo Segundo, artículos 323 a 329. \\

\textbf{¿Cuáles son los caracteres de las medidas preparatorias?} &
Son excepcionales o taxativas; no provocan la apertura de la instancia; interrumpen el curso de la prescripción; el juez competente es el que intervendrá en el proceso a promoverse; y pueden pedirse antes o después de interpuesta la demanda (siempre antes de la notificación del traslado). \\

\textbf{¿Qué significa que las medidas preparatorias se dicten ``sin sustanciación''?} &
Que se resuelven sin controvertir el trámite con la otra parte, no que carezcan de trámite. \\

\textbf{¿Qué sanciones prevé el artículo 329 ante el incumplimiento?} &
Multa, astreintes, presunción en contra del omiso, secuestro con eventual allanamiento, según la medida de que se trate. \\

\textbf{¿Cuándo caducan las medidas preparatorias?} &
A los 30 días de que se tuvo por cumplida la resolución que las ordena, si no se promueve la demanda, salvo las de los incisos 9, 10 y 11 del artículo 323. \\

\textbf{¿Por qué la prueba anticipada tiene naturaleza cautelar?} &
Porque busca evitar el perjuicio que produciría la pérdida de la prueba por el paso del tiempo, compartiendo con las medidas cautelares la finalidad de actuar ante el peligro en la demora. \\

\textbf{¿Qué diferencia a la prueba anticipada de las medidas preparatorias en cuanto a la contraparte?} &
En la prueba anticipada debe citarse a la contraparte (o al defensor oficial) para que controle la producción de la prueba, bajo pena de nulidad. \\

\textbf{¿Es apelable la denegatoria de la prueba anticipada?} &
Sí, como excepción al principio general de inapelabilidad de las resoluciones sobre prueba (artículo 379). \\

\midrule

\multicolumn{2}{p{\textwidth}}{
\textbf{Resumen:} Las diligencias preliminares comprenden dos institutos distintos: las medidas preparatorias (art. 323 a 329), de carácter taxativo, que no abren la instancia pero interrumpen la prescripción y caducan a los 30 días de cumplidas si no se promueve la demanda; y la prueba anticipada (art. 326), de naturaleza cautelar, que permite producir prueba antes de la etapa probatoria cuando existe riesgo de pérdida por el paso del tiempo, exigiendo la citación de la contraparte o del defensor oficial y admitiendo apelación en caso de denegatoria.
} \\

\bottomrule
\end{tabularx}
\end{table}

\chapter{La demanda}

\section{La acción, la pretensión y la demanda}

La acción es la posibilidad que tienen todos los ciudadanos de peticionar a las autoridades, conforme al artículo 14 de la Constitución Nacional, que excede al derecho procesal en tanto puede peticionarse en distintos ámbitos —ante el Poder Ejecutivo o ante el Poder Legislativo—. En lo que atañe específicamente al proceso judicial, la acción es el poder de hacer valer una pretensión ante la jurisdicción.

Una vez suscitado un conflicto —un choque de intereses entre determinados sujetos con posibilidad de trascender jurídicamente—, corresponde llevar ese conflicto a la sede de un tribunal, en tanto es el Poder Judicial de la Nación quien tiene la potestad, en los casos de su competencia, de decir el derecho y hacer actuar la voluntad de la ley.

\begin{definition}{Demanda (retomada)}
La demanda es un acto de petición que contiene una pretensión. Constituye, siguiendo este concepto, un acto de postulación, que está en cabeza de la parte juntamente con el letrado patrocinante, dado que, por el artículo 56, no pueden ejercerse actos postulatorios sin asistencia letrada.
\end{definition}

\section{La demanda como acto de postulación y sus efectos}

La demanda, como acto de postulación —junto con otros actos que se dan gracias a la acción y a los recursos—, transmite la apertura de la instancia, que constituye el ámbito adecuado para ejercer las demás peticiones que correspondan a lo largo del proceso.

\begin{important}
Una vez producido el efecto procesal de apertura de la instancia, quien inició el proceso debe efectuar los actos procesales correspondientes para hacerlo avanzar hacia el dictado de la sentencia. La caducidad de la instancia es un instituto que puede hacer perder los efectos producidos por la interposición de la demanda: si la instancia caduca, la demanda se tiene por no interpuesta, con consecuencias graves en materia de costas y, eventualmente, en la responsabilidad formal y hasta real del abogado que permitió que la instancia caducara.
\end{important}

\section{Clases de demanda}

Puede distinguirse entre demanda principal y demanda incidental.

% TODO: contenido incompleto o ambiguo en las notas originales — el apunte enuncia esta clasificación sin desarrollar las características particulares de cada clase.

\section{Acumulación objetiva de pretensiones (artículo 87)}

La acumulación objetiva de pretensiones está regulada en el artículo 87 del Código Procesal. Para su procedencia se exige que las pretensiones acumuladas no sean contradictorias entre sí, que puedan sustanciarse por los mismos trámites y que resulten de la competencia del mismo juez.

\section{Demanda mínima, admisible y fundada}

Esta clasificación permite ingresar en los requisitos exigidos por el artículo 330 para la interposición de una demanda.

\begin{definition}{Demanda mínima}
Debe contar, como mínimo, con una petición por escrito y firmada. Basta para activar la jurisdicción, en tanto el artículo 50 del Reglamento de la Justicia Nacional aclara que ningún escrito puede rechazarse de plano sin la resolución correspondiente del tribunal. Sin embargo, si la demanda se limita a esa petición mínima, será rechazada \textit{in limine} en los términos del artículo 337 del Código.
\end{definition}

\begin{definition}{Demanda admisible}
Además de indicarse con precisión su objeto, de estar hecha por escrito y de estar firmada —requisitos de la demanda mínima—, la demanda admisible debe cumplimentar los requisitos del artículo 330 del Código Procesal.
\end{definition}

\begin{definition}{Demanda fundada}
Puede ocurrir que se respeten los requisitos de admisibilidad de la demanda pero que la pretensión tenga un objeto ilícito —por ejemplo, si se pretende hacer valer un contrato de esclavitud o un contrato mediante el cual las partes acordaron participar de algún modo en la comisión de un delito—. En tal caso, la demanda puede ser admisible en cuanto a las formas, pero no fundada, por lo que será rechazada por improponibilidad objetiva, conforme al artículo 337, desde el inicio del proceso por parte del juez.
\end{definition}

\chapter{Requisitos de la demanda (artículo 330 CPCCN)}

El artículo 330 establece que la demanda se presentará por escrito. Ello se debe a que el proceso civil y comercial es preponderantemente escrito —no exclusivamente, dado que existen actuaciones orales—, aunque la gran mayoría de los actos constitutivos del proceso (demanda, reconvención, contestación de la reconvención, excepciones previas y contestación de excepciones previas), que integran la etapa introductoria del proceso, se realizan por escrito.

\section{Inciso 1: nombre y domicilio del demandante}

Debe individualizarse si la persona que se presenta lo hace por su propio derecho —cuando tiene capacidad procesal para ejercer por sí actos procesales válidamente— o si delega la representación convencional en otra persona.

\begin{itemize}
    \item Si se trata de una persona humana, corresponde indicar su nombre y apellido. El Código no exige el DNI, aunque resulta conveniente, casi indispensable, consignarlo, dado que pueden existir homónimos que generen confusión. El Código, en cambio, sí exige el nombre y el domicilio del demandante.
    \item Si se trata de una persona jurídica, debe indicarse su denominación social o razón social, además del domicilio correspondiente.
\end{itemize}

Respecto del domicilio: si se trata de una persona humana, corresponde denunciar el domicilio real, entendido como aquel en el que tiene su residencia habitual, conforme al artículo 73 del Código Civil y Comercial. Si se trata de una persona jurídica, corresponde denunciar el domicilio legal establecido por la ley 19.550, artículo 11.

\section{Inciso 2: nombre y domicilio del demandado}

Corresponden las mismas consideraciones efectuadas respecto del actor. Además de denunciar el domicilio real del demandado, corresponde constituir domicilio procesal, entendido como el domicilio del estudio del abogado (domicilio procesal físico), así como constituir domicilio procesal electrónico, que coincide con el número de CUIT del letrado.

Respecto del domicilio del demandado, cabe señalar que se trata de su domicilio real —si es persona humana— o legal —si es persona jurídica—, y que este dato suele conocerse recién cuando el demandado comparece a contestar la demanda. Para las personas jurídicas puede solicitarse un informe a la Inspección General de Justicia a fin de conocer el domicilio legal y evitar inconvenientes con la notificación del traslado.

\section{Inciso 3: la cosa demandada}

La cosa demandada es el objeto del litigio, es decir, la pretensión: aquello que se reclama que el juez otorgue en la sentencia definitiva. El Código exige que se designe con toda exactitud, exigencia que se corresponde con la obligación del juez de expedirse conforme al principio de congruencia.

\begin{important}
La pretensión delimita el poder del juez para decidir. En el proceso civil y comercial, que rige por el principio dispositivo, el juez solo interviene a instancia de parte y, por el principio de congruencia, se encuentra atado a las pretensiones de las partes.
\end{important}

Debe señalarse el objeto inmediato de la pretensión —qué tipo de sentencia se pide: constitutiva, declarativa o de condena— y el objeto mediato de la pretensión, que se corresponde con el bien de la vida que se reclama.

\section{Inciso 4: los hechos}

El inciso 4° del artículo 330 exige que se expliquen claramente los hechos en que se funda la demanda, lo cual implica explicar la causa de la pretensión, es decir, su fundamento de hecho.

Los hechos deben exponerse en forma clara, circunstanciada y en orden cronológico, en función del principio de sustanciación, teniendo en cuenta la norma que se invoca como fundamento de la pretensión, de manera tal que los hechos encuadren en el antecedente de aplicación de esa norma.

\begin{definition}{Hecho}
Un hecho es un acontecimiento o suceso que ocurre en el mundo exterior y que se percibe por la modificación que genera en la fuente de prueba, en las cosas o en las personas. Se conoce la existencia de un hecho en tanto algo se modifica en el mundo exterior; si no hay modificación alguna, o no puede inferirse la existencia del hecho por su fijación en otro, no puede concluirse que el hecho existe.
\end{definition}

El hecho, apenas sucede, queda en el pasado; para reconstruirlo en el proceso es necesario, en primer lugar, afirmarlo. En el capítulo de hechos de la demanda, para cumplimentar el inciso 4° del artículo 330, corresponde afirmar los hechos en forma cronológica, mediante oraciones de tipo narrativo, incorporando descripciones cuando corresponda, en sentido afirmativo, procurando claridad y evitando complicaciones innecesarias en el uso del lenguaje.

\begin{note}
Al redactar el capítulo de hechos debe observarse especial cuidado para relatarlos en forma cronológica y ordenada, sin dar lugar a confusión. Puede describirse, a modo de fotografía, la situación inicial en la que se vinculan las partes; luego desarrollarse, a modo de película, cómo se relacionaron entre sí; y finalmente realizarse una síntesis que concluya en que los hechos se encuadran en la norma elegida, que otorga responsabilidad a quien corresponda y que, por lo tanto, debe ser condenado a resarcir.
\end{note}

\section{Inciso 5: el derecho}

El inciso 5° del artículo 330 permite exponer el derecho sucintamente, es decir, brevemente, evitando repeticiones innecesarias, en tanto es el juez quien conoce el derecho y quien aplicará, al dictar sentencia, la norma que corresponda según los hechos acreditados en el proceso mediante los medios de prueba correspondientes.

\begin{important}
Independientemente de que se haya invocado correctamente el derecho, o incluso de que se lo haya omitido, el juez suplirá esta circunstancia y aplicará la norma que corresponda a los hechos acreditados en la causa. Esta posibilidad de exponer el derecho sucintamente debe revisarse cuando se trate de una cuestión eminentemente jurídica —por ejemplo, cuando se discuta qué norma rige el caso o qué interpretación corresponde darle—, supuesto en el cual la carga de exponer adecuadamente las razones jurídicas que fundamentan la pretensión recae en la parte.
\end{important}

\section{Inciso 6: la petición en términos claros y positivos}

El inciso 6° del artículo 330 exige que la petición se efectúe en términos claros y positivos, en consonancia con la obligación del juez, en virtud del principio de congruencia, de expedirse en relación con las pretensiones y los hechos afirmados por las partes.

Debe precisarse el monto reclamado: no basta con indicar qué se reclama, sino que, si se reclama una suma de dinero, debe precisarse el monto correspondiente. Se exceptúa el caso en que, por determinadas circunstancias, no sea posible hacerlo al momento de interponer la demanda, o en que ello pudiera perjudicar el curso de la prescripción; en tal supuesto corresponde reservar la determinación exacta del monto para lo que resulte de la prueba producida en el expediente.

\section{La prueba (artículo 333): lo que se acompaña con la demanda}

El artículo 333 regula la prueba que debe incorporarse junto con la demanda. Conforme al régimen actual, deben incorporarse todos los medios de prueba de los que la parte pretenda valerse para acreditar los hechos afirmados como fundamento de su pretensión.

\chapter{Efectos de la demanda}

\section{Efectos procesales}

El principal efecto procesal de la interposición de la demanda es la apertura de la instancia, que significa la posibilidad de realizar actos procesales ante el juez.

\begin{important}
La apertura de la instancia tiene por finalidad, además, hacer comenzar a correr el plazo de caducidad de la instancia.
\end{important}

\section{Efectos sustanciales}

El principal efecto sustancial de la interposición de la demanda es que interrumpe el curso de la prescripción.

\section{Requisitos fiscales y mediación previa obligatoria}

Junto con los requisitos de admisibilidad, deben valorarse también los requisitos fiscales, en tanto la demanda constituye una acción altamente imponible, cuyo pago corresponde a quien la promueve, salvo que goce del beneficio de litigar sin gastos. La ley 3.898 establece el pago de la tasa de justicia, que debe efectuarse al inicio del proceso judicial.

% ---------------------------------------------------------
% CORNELL 2 — LA DEMANDA
% ---------------------------------------------------------
\section*{Repaso: la demanda y sus requisitos}
\addcontentsline{toc}{section}{Repaso: la demanda y sus requisitos}

\begin{table}[H]
\centering
\small
\renewcommand{\arraystretch}{1.05}
\setlength{\parskip}{0.2em}
\begin{tabularx}{\textwidth}{
    >{\raggedright\arraybackslash}p{0.30\textwidth}
    >{\raggedright\arraybackslash}X
}
\toprule
\textbf{Preguntas / palabras clave} &
\textbf{Notas / respuestas} \\
\midrule

\textbf{¿Qué es la demanda?} &
Un acto de petición que contiene una pretensión, y que constituye un acto de postulación en cabeza de la parte y del letrado patrocinante. \\

\textbf{¿Qué exige el artículo 87 para la acumulación objetiva de pretensiones?} &
Que las pretensiones no sean contradictorias entre sí, que puedan sustanciarse por los mismos trámites y que resulten de la competencia del mismo juez. \\

\textbf{¿Qué diferencia hay entre demanda mínima, admisible y fundada?} &
La mínima solo requiere petición por escrito y firmada; la admisible debe cumplir los requisitos del artículo 330; la fundada, además, debe tener un objeto lícito y una pretensión proponible. \\

\textbf{¿Qué debe indicarse en el inciso 1° del artículo 330?} &
El nombre y domicilio del demandante (real si es persona humana, legal si es persona jurídica). \\

\textbf{¿Qué domicilios deben constituirse respecto del demandado?} &
El domicilio real (o legal) y, además, el domicilio procesal (físico, en el estudio del abogado) y el domicilio procesal electrónico (CUIT del letrado). \\

\textbf{¿Qué es la cosa demandada (inciso 3°)?} &
El objeto del litigio, es decir, la pretensión que se reclama al juez que otorgue en la sentencia, comprendiendo objeto inmediato y objeto mediato. \\

\textbf{¿Cómo deben exponerse los hechos (inciso 4°)?} &
En forma clara, circunstanciada y en orden cronológico, conforme al principio de sustanciación. \\

\textbf{¿Cómo debe exponerse el derecho (inciso 5°)?} &
Sucintamente, evitando repeticiones innecesarias, ya que el juez conoce el derecho y lo suple si fuera necesario. \\

\textbf{¿Qué exige el inciso 6° respecto de la petición?} &
Que se formule en términos claros y positivos, precisando el monto reclamado cuando corresponda. \\

\textbf{¿Qué regula el artículo 333?} &
La prueba que debe incorporarse junto con la demanda, ofreciendo todos los medios de que la parte quiera valerse. \\

\textbf{¿Cuáles son los efectos procesales y sustanciales de la demanda?} &
El efecto procesal principal es la apertura de la instancia (que además hace correr el plazo de caducidad); el efecto sustancial principal es la interrupción de la prescripción. \\

\midrule

\multicolumn{2}{p{\textwidth}}{
\textbf{Resumen:} La demanda es el acto de postulación mediante el cual se ejerce la acción y se introduce una pretensión ante el juez. Debe cumplir los requisitos del artículo 330 —individualización de las partes y sus domicilios, la cosa demandada, los hechos, el derecho y una petición clara y positiva— y acompañarse de la prueba conforme al artículo 333. Su interposición produce la apertura de la instancia como efecto procesal y la interrupción de la prescripción como efecto sustancial, además de generar obligaciones fiscales para quien la promueve.
} \\

\bottomrule
\end{tabularx}
\end{table}

\end{document}
